% =============================================================================
% Deep Learning for Parameter Estimation in Spatial Point Processes
% RESTRUCTURING PASS -- no prose has been rewritten, only moved.
%
% HOW TO READ THIS FILE
%   * Every section opens with a % NOTE block stating its PURPOSE and its TODOs.
%   * \todo{...} marks inline gaps (renders red while \drafttrue).
%   * Text carried over from the old draft is tagged % [MOVED FROM: ...].
%   * Sections 2, 7.1, 7.3, 8, 9 are currently empty skeletons -- these are the
%     writing targets. Sections 5 and 6 are largely done.
% =============================================================================

\documentclass[11pt,letterpaper]{article}

\usepackage[margin=1in]{geometry}
\usepackage{mathpazo}

\usepackage{amsmath}
\usepackage{amssymb}
\usepackage{amsthm}
\usepackage{mathtools}
\usepackage{dsfont}

\usepackage{graphicx}
\usepackage{booktabs}
\usepackage{enumitem}

\usepackage{natbib}
\usepackage{xcolor}
\usepackage[hidelinks]{hyperref}

% ----------------------------------------------------------------------------
% Macros (unused ones from the AAAI draft have been deleted)
% ----------------------------------------------------------------------------
\newcommand{\R}{\mathbb{R}}
\newcommand{\N}{\mathbb{N}}
\newcommand{\Z}{\mathbb{Z}}
\newcommand{\E}{\mathbb{E}}
\newcommand{\1}{\mathds{1}}
\newcommand{\pc}{\mathbf{x}}
\newcommand{\simpcom}{\Sigma}
\newcommand{\filt}{\mathcal{F}}
\newcommand{\kpar}{\kappa}
\newcommand{\mo}{\mu}

% ----------------------------------------------------------------------------
% Draft annotations -- set \draftfalse to hide all TODOs
% ----------------------------------------------------------------------------
\newif\ifdraft
\drafttrue
\ifdraft
  \newcommand{\todo}[1]{\textcolor{red}{[TODO: #1]}}
  \newcommand{\fix}[1]{\textcolor{blue}{[FIX: #1]}}
\else
  \newcommand{\todo}[1]{}
  \newcommand{\fix}[1]{}
\fi

% ----------------------------------------------------------------------------
% Theorem environments (trimmed to those actually used)
% ----------------------------------------------------------------------------
\theoremstyle{plain}
\newtheorem{theorem}{Theorem}[section]
\newtheorem{claim}[theorem]{Claim}

\newtheoremstyle{definitioncolon}%
  {\topsep}{\topsep}{\normalfont}{}{\bfseries}{}{.5em}%
  {\thmname{#1}\thmnumber{ #2}\thmnote{: #3}}
\theoremstyle{definitioncolon}
\newtheorem{definition}[theorem]{Definition}

\newtheorem*{note}{Note}
\newtheorem*{example}{Example}
\newtheorem*{remark}{Remark}

% ----------------------------------------------------------------------------
\title{Topological Feature Learning for Parameter Estimation\\
       in Neyman--Scott Processes
       \thanks{\todo{title is a placeholder -- it should name the topological
       angle and the family class, not just ``deep learning''}}}
\author{Flavio Gualtieri}
\date{\today}

\begin{document}
\maketitle


% #############################################################################
% ABSTRACT
% #############################################################################
% NOTE -- PURPOSE: sell the paper in 6 sentences.
% The current version is an inventory of components (pipeline, encoder, FCNN).
% It should instead be: problem -> limitation of classical estimators ->
% topological features as the answer -> TWO threads (which summary is best /
% how good is the best model) -> multi-family coverage -> headline number.
% REWRITE LAST, once Sections 1 and 7 are settled.
% #############################################################################
\begin{abstract}
% [MOVED FROM: old abstract -- keep as raw material only, do not ship]
Overview of deep learning estimator for the parameters of homogeneous spatial
point processes. It is comprised of (1) a persistence vectorization pipeline,
(2) an encoder that produces a feature vector, and (3) a FCNN with $p$ output
nodes corresponding to the $n$ parameters of the SPP.

\todo{rewrite entirely; leave a bracketed slot for the headline result}

Problem: why classical estimators are limited
We propose a method using vectorized persistence summaries as learned features, 
and explore (i) which summaries are most effective, and (ii) how the best model 
compares to exsting baselines.
\end{abstract}


% #############################################################################
% 1. INTRODUCTION
% #############################################################################
% NOTE -- PURPOSE: motivate, position, and promise. Four moves in this order:
%   (i)   parameters are the interpretable scientific object
%   (ii)  classical estimators depend on a CHOSEN summary statistic with a
%         closed form -- state each limitation as a claim, not a noun
%   (iii) topological summaries: rich, family-agnostic, no closed form needed
%   (iv)  explicit contributions + roadmap
%
% BIGGEST PROBLEM: the word "topology" does not currently appear anywhere in
% this section. The reader hits Section 3.2 with no idea why simplicial
% complexes are being introduced. Fix (iii) first.
% #############################################################################
\section{Introduction}

% --- (i) motivation ----------------------------------------------------------
% [MOVED FROM: old Section 1]
The study of spatial point processes finds applications in astronomy, cell
biology, epidemiology, etc. Many natural phenomena - such as galaxy clustering -
are modelled using SPPs, so one can hope to leverage an understanding of these
processes to obtain novel insights.
\todo{pick ONE application and carry it through the whole paper as the running
example -- galaxy clustering is the natural choice given the related work}

The ability to infer the parameters of SPPs is of particular interest because a
wealth of canonical summary statistics become readily available once these
parameters are known, such as ...
\todo{this motivation is CIRCULAR and must be replaced. Classical methods
estimate parameters FROM summary statistics, not the other way round. The real
argument: the parameters are the physically interpretable quantities
($\kappa$ = cluster density, $\mu$ = occupancy, $\sigma$ = cluster scale), and
existing estimators are fragile in nameable regimes.}
\fix{The parameters of spatial point processes are of particular interest because 
they are interpretable scientific quantities. For example ...}


% --- (ii) limitations of existing estimators ---------------------------------
A number of parameter estimation methods have been proposed, typically relying
on closed-form expressions for summary statistics, for example the $K$ and $g$
functions.
\todo{these three become a short paragraph with citations, not a bullet list}
\begin{itemize}[nosep]
  \item Minimum contrast on $K$ or $g$ (Diggle 2003).
  \item Palm likelihood (Tanaka, Ogata \& Stoyan 2008).
  \item Bayesian and MCMC.
\end{itemize}

These methods have limitations we hope to address.
\todo{THESE FOUR BARE NOUNS ARE THE LOAD-BEARING JUSTIFICATION FOR THE ENTIRE
PROJECT and are currently placeholders. Each must become a claim with a
mechanism and, ideally, a regime where it bites:
 - hand-picked statistic: $K$ radially averages away multi-scale structure;
   a closed form may not exist at all (nested Thomas)
 - hyperparameter sensitivity: minimum contrast depends on $(r_{max}, q, p)$
 - stability: which regimes? near-Poisson? name them
 - generalizability: new family => new derivation}
\begin{itemize}[nosep]
  \item Reliance on hand-picked summary statistic (and by extension a
        closed-form expression for this statistic).
  \item Hyperparameter tuning sensitivity.
  \item Stability.
  \item Generalizability.
\end{itemize}
\fix{Firstly, existing methods rely on the existence of closed-form expressions of
summary statistics -- such as the $K$-function -- which may not exist for every process.
For example, while such an expression of the $K$-function exists for the Thomas process,
it does not for the nested Thomas process. Secondly, classical parameter estimation
methods suffer from hyperparameter sensitivity, specifically the minimum contrast method
depends on $(r_{max}, q, p)$. Thirdly, classical methods are not stable in nameable regimes.
Finally, these methods do not generalize well to new families of processes, as a new derivation
is required for each new family.}

% --- (iii) the topological turn ---------------------------------------------
\todo{MISSING PARAGRAPH. This is the hinge of the paper: persistent homology
gives a multi-scale, family-agnostic descriptor of the pattern that requires no
closed form and no choice of statistic. Without this paragraph, Section 3.2
appears out of nowhere.}
\fix{These limitations motivate the use of persistence-based features which provide a
multi-scale, process-agnostic descriptor of the point cloud that does not require
a closed-form or a choice of summary statistic.}

% --- (iv) contributions ------------------------------------------------------
\todo{MISSING CONTRIBUTIONS LIST. Draft it as:
 1. a systematic comparison of vectorized topological summaries as features for
    SPP parameter estimation (the ORIGINAL objective -- re-elevate it);
 2. a per-diagram coverage-quantile calibration scheme for vectorizing across a
    heterogeneous population of diagrams, with an explicit leakage protocol;
 3. a validated multi-family simulator + evaluation protocol, including an exact
    finite-window variance identity;
 4. an estimator covering Thomas / Mat\'ern cluster / nested Thomas under a
    common reparameterization.
NOTE: architectural novelty is NOT on this list -- see Section 2.}

\todo{roadmap paragraph}


% #############################################################################
% 2. RELATED WORK
% #############################################################################
% NOTE -- PURPOSE: locate the work honestly. Target ~half a page. ENTIRELY NEW.
%
% This section exists because the "novel architecture" framing is dead. Two
% works occupy the neighbourhood:
%   [A] "Using neural networks to estimate parameters in spatial point process
%       models" -- essentially this architecture, but with the L function as
%       the feature.
%   [B] "Learning from topology: cosmological parameter estimation from the
%       large-scale structure" -- topological features, but for cosmological
%       parameters.
%
% STRUCTURE: position on two axes -- (what the feature is) x (what is being
% estimated). [A] establishes the architecture pattern. [B] establishes
% topological features for parameter inference. Concede both plainly and early;
% the paper is stronger for it.
%
% THEN state the gap precisely. Candidates, in descending order of defensibility:
%   - neither systematically COMPARES topological summaries against each other
%     as estimator features (this is the thread to own)
%   - neither addresses calibration across a heterogeneous diagram population
%   - neither covers multiple Neyman--Scott families under one reparameterization
%   - [B] is cosmology-specific; the SPP-estimation setting has a known error
%     floor and classical baselines to measure against
% #############################################################################
\section{Related Work}

\todo{WRITE THIS SECTION. Concede the architecture to [A]; concede topological
features for parameter estimation to [B]; claim the comparative study, the
calibration methodology, and the multi-family scope.}

\todo{one paragraph: neural estimators for SPPs -- [A] and its lineage}
\fix{REF first proposed the use of deep learning for parameter estimation. They implement
an encoder/inference architecture similar to PointNet, but use the $L$-function as the learned
feature and therefore suffer from the same limirations as classical methods (i.e. reliance
on a closed-form summary statistic). REF employ a similar architecture as REF, but actually use
topological features. Their model is deigned to estimate cosmological parameters rather than the
parameters of the underlying point process. Furhtermore, they do not systematically compare 
various topological summaries, nor do they address the calibration of the summaries they use.}

\todo{one paragraph: topological summaries as features for inference -- [B],
plus the vectorization literature (persistence images, Betti curves)}
\fix{The vectorization of persistence diagrams is an active area of research arising from the issue
that persistence diagrams -- not being a Hilbert space -- cannot be direclty used as inputs for
machine learning models. To mitigate this, REF propose persistence images, which vectorise persistence
diagrams by applying a Gaussian kernel to each point in the diagram and intergating over a grid.
REF propose Betti curves, which vectorize persistence diagrams by counting the number of topological features
that are alive at each scale.}

\todo{one paragraph: the gap, stated in one or two sentences}


% #############################################################################
% 3. BACKGROUND
% #############################################################################
% NOTE -- PURPOSE: give the reader only what is needed downstream.
% Currently ~40% of the draft and written as neutral textbook exposition.
%
% ACTIONS:
%   * CUT hard. Keep intensity, Neyman--Scott, the three families; keep
%     filtration -> diagram -> vectorization.
%   * MOVED OUT: the $g$ / $K$ definitions now live in Section 5.3, which is
%     the only place they are used.
%   * ADD the missing justification paragraph at 3.2 (see \todo there).
%   * The literature-review placeholder lives at the end of this section.
% #############################################################################
\section{Background}
\label{sec:background}

% [MOVED FROM: old Section 2 preamble]
We provide the background in SPPs and persistent homology required for this
discussion. Readers with a working understanding of both may skip to
Section \ref{sec:method}.

% -----------------------------------------------------------------------------
\subsection{Spatial point processes}
% NOTE: keep as is. This is tight and load-bearing.
% -----------------------------------------------------------------------------

    \begin{definition}[Intensity function]
        A $d$-dimensional SPP $X$ on a subset $W \subseteq \R^d$ is equipped
        with an intensity function $\rho(u): \R^d \rightarrow \R^+$ such that
            \[
            \E[N(W)] = \int_W \rho(u) du,
            \]
        where $N(W) = \mid W \cap X \mid$ denotes the number of points of $X$
        in $W$. When $\rho(u) = \lambda \in \R$, $X$ is said to be
        \emph{homogeneous}.
    \end{definition}

    \begin{example}
        For a homogeneous Poisson process with intensity $\lambda$, we have
        $\rho(u) = \frac{1}{\lambda}$, i.e. points are scattered uniformly on
        the space.
    \end{example}

    \begin{definition}[Neyman-Scott process]
        Let $P$ be a homogeneous Poisson \emph{parent} process with intensity
        $\kpar$. For a Neyman-Scott process with Poisson-distributed offspring
        we have:
            \[
            \rho_{NS}(u) = \mo \sum_{p \in P} k(u - p),
            \]
        where $k$ is a kernel that determines how the \emph{offspring} points
        are distributed around the parents.
    \end{definition}

    \begin{note}
        Neyman-Scott processes always have the parent and offspring intensities
        as parameters. Additional parameters are introduced by the choice of
        kernel.
    \end{note}

    \begin{definition}[Homogeneous Thomas process]
        A homogeneous Thomas process is a Neyman-Scott process with an isotropic
        Gaussian kernel $\mathcal{N}(0, \sigma^2I)$
            \[
            \rho_{Thomas}(u) = \mo \sum_{p \in P} k(0, \sigma^s),
            \]
        This process has parameters $\theta_{Thomas} = (\kappa, \mu, \sigma)$
        corresponding to the parent intensity, offspring intensity, and cluster
        scale, respectively.
    \end{definition}

\todo{add the Mat\'ern cluster and nested Thomas definitions here -- both are
used in Section 5 but never defined}

% -----------------------------------------------------------------------------
\subsection{Persistent homology}
\label{ssec:persistence}
% NOTE: mechanically complete but intellectually silent. It defines everything
% EXCEPT the one thing the reader needs.
% -----------------------------------------------------------------------------

\todo{OPENING PARAGRAPH MISSING -- and it is the intellectual justification for
the whole method. What does persistence see in a clustered pattern that a
radially averaged second-order statistic does not? Multi-scale connectivity
($H_0$ death times track cluster merging across scales) and void structure
($H_1$). Write this BEFORE the definitions, so the reader knows what the
machinery is for.}
\fix{Redially averaged statistics -- such as the $g$ and $K$ functions -- collapse
a point cloud to a real-valued curve indexed by scale $r$. Features based on such
statistics are constructed by evaluating over a chosen range $[r_{min}, r_{max}]$
and subsampling at $n$ points along the interval to obtain a feature vector in $\R^n$.
This approach is limited because it discards multi-scale connectivity information and
void structure, i.e. \emph{how} the density is distributed across the space.
Persistent homology retains this information by tracking a family of simplicial
complexes on the point cloud as the scale parameter grows, instead of a single scalar
summary. Under a filtration such as Vietoris-Rips, the $H_0$ persistence summary records
the order and scales at which clusters merge. The $H_1$ counterpart marks voids and gaps
in the pattern as they persist across scales, allowing one to capture the clustering-induced
emptiness or repulsion that a process may exhibit, which a radial-average cannot see.
For each $d \in \N$, the $d$-dimensional persistence summary requires no closed form expression
and no assumption about which process family generated it.
}

% [MOVED FROM: old Section 2.3 -- unchanged]
Let $X$ be a point process on $R^d$ and $W \subseteq \R^d$ a bounded window with
$|X \cap W| = n$ . Enumerate the points of $X \cap W$ to construct a point cloud
$\pc = \{x_1, ..., x_n\} \subset \R^d$. Let $\Sigma(\pc)$ denote the \emph{full
simplex} on $\pc$, i.e. the collection of all non-empty subsets of $\pc$

    \begin{definition}[Simplex]
        A \emph{simplex} of $\pc$ is a non-empty finite subset
        $\sigma \subseteq \pc$. Its faces are the non-empty subsets
        $\tau \subseteq \sigma$.
    \end{definition}

    \begin{definition}[Simplicial complex]
        A simplicial complex $\simpcom$ on $\pc$ is a collection of simplices
        closed under taking faces: $\sigma \in \simpcom, \tau \subseteq \sigma
        \implies \tau \in \simpcom$.
    \end{definition}

    \begin{definition}
        A filtration function on $\pc$ is a monotone map
        $f: \Sigma(\pc) \rightarrow \R^+$ such that
            \[
            f(\tau) \leq f(\sigma) \quad \forall \quad \tau \subseteq \sigma.
            \]
    \end{definition}

    \begin{example}
        Vietoris-Rips filtration function:
            \[
            f_{VR}(\{x_i\}) = 0, \qquad f_{VR}(\{x_i, x_j\}) = \|x_i - x_j\|,
            \]
        extended to higher simplices by $f_{VR}(\sigma) = \max_{\{x_i,x_j\}
        \subseteq \sigma} f_{VR}(\{x_i, x_j\})$, i.e. the diameter of $\sigma$.
    \end{example}

    \begin{example}
    \label{ex:dtm}
        For $u \in \R^d$, let $N_k(u)$ denote its $k$ nearest neighbours in
        $\pc$. The $\mathrm{DTM}_k$ filtration function is
            \[
            f_{DTM}(\sigma; k) = \max_{u \in \sigma}\{(\frac{1}{k}
            \sum_{v \in N_k(u)} d(u, v)^2)^{\frac{1}{2}}\}
            \]
        again extended to higher simplices by the max-over-edges rule of the VR
        example. \citep{anai2020dtm}.
    \end{example}

    \begin{claim}
        $f_{DTM_\kappa}$ is monotone, and is hence a valid filtration function
        \citep[Prop.\ 3.5]{anai2020dtm}.
    \end{claim}

    \begin{definition}
        Given a filtration function $f$ on $\pc$ and $r \geq 0$, we construct
            \[
            \mathcal{F}_r = \{\sigma \in \Sigma(\pc) : f(\sigma)
            \leq r\}.
            \]
        Monotonicity of $f$ ensures each $\mathcal{F}_r$ is a simplicial complex
        and that $\mathcal{F}_s \subseteq \mathcal{F}_t$ for $0 \leq s \leq t$.
        The family $(\mathcal{F}_r)_{r \geq 0}$ is the \emph{filtration} induced
        by $f$. We write $VR_r(\pc)$ for $\mathcal{F}_r$ under $f_{VR}$.
    \end{definition}

    \begin{definition}[Betti number]
        Fix $k \in \N$ and a simplicial complex $\Sigma$. Taking homology
        with coefficients in a field (we use $\Z_2$), $H_k(\simpcom)$ is a
        vector space of dimension $\beta_k = \dim H_k(\simpcom)$, the $k$-th
        \emph{Betti number}. $\beta_0, \beta_1, \beta_2, \dots$ count connected
        components, loops, and cavities, respectively.
    \end{definition}

    \begin{definition}[Persistence pair \& lifetime]
        For a filtration $(\filt_r)_{r \geq 0}$, a $k$-dimensional homology
        class is \emph{born} at the smallest scale $b \geq 0$ at which it
        appears in $H_k(\simpcom_b)$, and \emph{dies} at the scale $d > b$ at
        which it merges into an older class or becomes a boundary. The pair
        $(b,d)$ is a \emph{persistence pair}, and $d - b$ its \emph{lifetime}.
    \end{definition}

    \begin{example}
        Consider $H_0$. Every $0$-dimensional class is born at $r=0$ and dies
        when its two components merge, i.e. at the filtration value of the
        connecting edge.
    \end{example}

    \begin{definition}[Persistence diagram]
        A persistence diagram $D_k = \{(b_i, d_i)\}_i$ is the multiset of
        dimension-$k$ persistence pairs, plotted with birth on the $x$-axis and
        death on the $y$-axis.
    \end{definition}

    \begin{remark}
        Every point of $D_k$ lies above the diagonal $b=d$ because it must die
        after it is born. Distance from the diagonal is lifetime, so short-lived
        points near the diagonal are often treated as noise.
    \end{remark}

    \begin{definition}[Persistence image]
    \label{def:per_image}
        The persistence image of $D_k$ is constructed by:
            \begin{enumerate}
                \item Mapping to birth-persistence coordinates:
                      $(b,d) \mapsto (b, d-b)$.
                \item Weighting each point by $w(d-b)$.
                \item Forming the persistence surface
                      $\rho(z) = \sum_i w(d_i - b_i)\,
                      \phi(z; (b_i, d_i - b_i))$, a weighted sum of isotropic
                      Gaussian kernels with bandwidth $\sigma_{im}$ centered at
                      the transformed points.
                \item Integrating $\rho$ over a $64 \times 64$ pixel grid,
                      giving a single-channel image $PI \in \R^{64 \times 64}$.
            \end{enumerate}
    \end{definition}

    \begin{definition}[Betti curve]
        For fixed homology dimension $k$, the Betti curve is the
        right-continuous step function counting $k$-features alive at scale $r$:
            \[
            \beta_k(r) = \#\{i : b_i \leq r < d_i\}
            = \sum_i \1[b_i \leq r < d_i].
            \]
    \end{definition}

\todo{list here the OTHER vectorizations that were implemented and tested --
they are the objects compared in Section 7.1 and must be defined somewhere}

% -----------------------------------------------------------------------------
\subsection{Background reading}
% NOTE: PROGRESSION-REPORT SLOT. Deliberately deferred to last.
% For the report this doubles as the literature-review requirement.
% -----------------------------------------------------------------------------
\todo{WRITE LAST. Reading clusters to cover:
 (a) classical SPP estimation -- Diggle, minimum contrast, Palm likelihood,
     composite likelihood;
 (b) TDA foundations -- persistence, stability, DTM filtrations;
 (c) vectorization of persistence -- persistence images, landscapes, Betti
     curves, kernel methods;
 (d) neural/simulation-based inference -- amortized inference, neural
     estimators, the two related works of Section 2;
 (e) TDA applied to point processes and to cosmology.}


% #############################################################################
% 4. PROBLEM SETUP
% #############################################################################
% NOTE -- PURPOSE: state formally what is being estimated and what "good" means,
% BEFORE any architecture is described. NEW SECTION, assembled from material
% that was previously scattered across the top of the old Model section and the
% bottom of data generation.
% Short: half a page.
% #############################################################################
\section{Problem Setup}
\label{sec:problem}

\todo{opening: formal statement of the estimation target -- given a single
realization $\pc$ on window $W$ from a Neyman--Scott process with unknown
$\theta$, produce $\hat\theta$. Emphasise SINGLE realization: this is what makes
the problem hard and it is what sets the error floor.}
\fix{Let $X$ be a Neyman-Scott process on $\R^d$ with unknown parameter vector
$\theta \in \Theta$. For example, the homogeneous Thomas process has 
$\theta = (\kappa, \mu, \sigma)$. We observe a single realization of $X$ on a 
bounded window $W \subset \R^d$, giving a point cloud
$\pc = {x_1, \dots, x_n} \subset \R^d$ with $n = |X \cap W|$. The objective is to 
produce an estimate $\hat\theta = \hat\theta(\pc)$ of $\theta$ from $\pc$ alone.
We stress \emph{single}: $\theta$ is not recovered from repeated draws of $X$, nor 
is $\pc$ pooled or averaged across realizations. This is the setting a practitioner
faces: one observed pattern and one chance to estimate the generating parameters.
It is also what makes the problem hard in a way no estimator can fix: a 
fixed $\theta$ already induces a distribution over point clouds, so even the true 
parameter value gives rise to a spread of plausible $\pc$, and it is this irreducible 
spread —- not any weakness of a particular estimator -— that sets the error floor 
characterized below.}


% [MOVED FROM: old Section 3, loss + log-normalization]
The model is trained on a training set of synthetic clouds
$\{\pc_i, \theta_i\}_{i=1}^{N_{train}}$ by backpropagating the RMS loss
    \[
    \mathcal{L}(\hat\theta_i, \theta_i) = \sqrt{\sum_{j=1}^{d}(\hat\theta^j_i - \theta^j_i)^2},
    \]
where $\theta_i = (\theta_i^1, \dots, \theta_i^d)$. Since parameters on a larger scale
dominate this loss, we train in log-normalized parameter space,
    \[
    \{\pc_i, \theta_i\}_{i=1}^{N_{train}} \rightarrow
    \left\{\pc_i, \frac{\log(\theta_i) - \mu_{\log\theta}}
    {\sigma_{\log\theta}}\right\}_{i=1}^{N_{train}},
    \]
where -- abusing notation -- $\mu_{\log\theta}, \sigma_{\log\theta}$ denote the
mean and standard deviation of the logged parameters.

% [MOVED FROM: end of old Section 3.1.1 -- belongs here, as the definition of
%  what success means]
We draw one realization per parameter vector. The dispersion of $\hat\theta$
across realizations sharing a common $\theta$ gives the irreducible error floor
against which the model's performance should be read.

\todo{make the error floor a named quantity here and reference it from
Section 7.3 -- it is the yardstick for every reported number}


% #############################################################################
% 5. SIMULATOR AND EXPERIMENTAL DESIGN
% #############################################################################
% NOTE -- PURPOSE: this was buried as a subsection of "Model". PROMOTED.
% Frame it as THE BENCHMARK, not as setup: it is a contribution in its own
% right (validated multi-family simulator + exact finite-window identity +
% evaluation protocol). It is also the strongest evidence of rigour available
% for the progression report -- it must be visible in the table of contents.
%
% This section is essentially DONE. Light edits only.
% #############################################################################
\section{Simulator and Experimental Design}
\label{sec:simulator}

% [MOVED FROM: old Section 3.1]
Training data consists of pairs $(\pc_i, \theta_i)$ in which $\theta_i$ is drawn
from a design distribution $\Pi$ over parameters space and each $\pc_i$ is a
single realization of the corresponding process. We describe the design
distribution, sampling algorithm, and verification procedure used to generate
the synthetic data.

% -----------------------------------------------------------------------------
\subsection{Design distribution}
% NOTE: done. The reparameterization argument is good and should be kept intact.
% -----------------------------------------------------------------------------
The naive approach of sampling $(\kappa, \mu, \sigma)$ uniformly and
independently from intervals is inadequate. Firstly, point processes can
typically be separated into qualitative regimes of interest, which occupy curved
subsets of the parameter grid. Therefore uniform sampling would allocate most of
its mass to configurations that are either near-Poisson or degenerate. Secondly,
the raw parameters are not comparable across families, for example Mat\'ern's
ball radius $R$ and Thomas's cluster scale $\sigma$ do not play equivalent roles.
We therefore sample in reparameterized coordinates that carry the same meaning
for every family considered:
    \begin{align*}
        \lambda &= \kappa\mu && \text{expected offspring intensity,} \\
        \kappa & && \text{parent intensity,} \\
        \nu &= r\sqrt{\kappa} && \text{dimensionless overlap index,}
    \end{align*}
where $r$ denotes the cluster scale of the family in question: $r = \sigma$ for
Thomas, $r = R$ for Mat\'ern cluster, and $r = \sqrt{\sigma_1^2 + \sigma_2^2}$
for nested Thomas. The overlap index $\nu$ measures cluster radius against mean
inter-parent spacing, so $\nu \ll 1$ gives well-separated clusters and
$\nu \gtrsim 1$ gives clusters that merge (so the process approaches complete
spatial randomness). Nested Thomas carries an additional coordinate, the scale
ratio $\rho = \sigma_2/\sigma_1 \in (0,1)$, which separates configurations in
which both levels of structure are resolvable from those in which the process
collapses to a single scale.

We draw $(\log\kappa, \log\lambda, \log\nu)$ uniformly over a grid and invert to
recover $\mu = \lambda/\kappa$ and $r = \nu/\sqrt{\kappa}$. Log-uniform sampling
is used because the parameters span several orders of magnitude, and it matches
the log-normalization applied to the ground-truth targets (see
Section \ref{sec:problem}) so that training targets are approximately uniform in
each coordinate. Draws are taken from a scrambled Sobol sequence.

The ranges of the design grid are constrained as follows:
    \begin{enumerate}
        \item \textbf{Point budget.} The cost of computing persistence summaries
        grows rapidly with the number of points, so we impose
        $\mathbb{E}[N] = \lambda \in [150, 800]$, fixing the range of
        $\log\lambda$.
        \item \textbf{Clusters per window.} We require $\kappa|W| = \kappa \geq 15$
        so that a single realization is informative about $\kappa$, and
        $\kappa|W| = \kappa \leq 120$ so that individual clusters remain
        resolvable.
        \item \textbf{Overlap.} We impose $\nu \in [0.10, 0.90]$. The upper
        limit deliberately admits the near-Poisson regime so that the model is
        still exposed to configurations in which the parameters are weakly
        identifiable.
    \end{enumerate}
Parameter draws violating these constraint are rejected and redrawn before any
cloud is generated. This filter acts on the design and not on realizations,
therefore it does not distort the conditional law of $\pc$ given $\theta$. The
realized acceptance rate is $24.8\%$.

% -----------------------------------------------------------------------------
\subsection{Sampling algorithm}
% NOTE: done.
% -----------------------------------------------------------------------------
The various instances of a the Neyman--Scott construction may be generated by a
single algorithm. The only part that is different for each instance is the
displacement kernel $k$ and the offspring count distribution. For a window
$W \subset \R^d$:
    \begin{enumerate}
        \item Expand $W$ by an edge buffer $b$ to produce $W_{exp}$, so that
        clusters whose parents lie outside $W$ are not truncated.
        \item Sample the parent process on $W_{exp}$:
        $n \sim \mathrm{Poisson}(\kappa|W_{exp}|)$, with parents placed uniformly
        on $W_{exp}$.
        \item For each parent $p$ independently, draw an offspring count
        $n_p \sim \mathrm{Poisson}(\mu)$.
        \item For each offspring independently, draw a displacement
        $\epsilon \sim k$ and place the point at $p + \epsilon$.
        \item Discard the parents and retain the offspring lying in $W$.
    \end{enumerate}
For the Thomas process $k = \mathcal{N}(0, \simpcom^2 I_d)$; for the Mat\'ern
cluster process $k$ is uniform on the ball of radius $R$; for nested Thomas the
parent process at step 2 is itself a Thomas process rather than a homogeneous
Poisson process.

The buffer is set by a single rule applied to whichever kernel is in use: $b$ is
the radius containing all but $\varepsilon$ of the mass of $k$. For a Gaussian
kernel this gives $b = 4\sigma$, at which the per-axis tail mass is
$2\Phi(-4) \approx 6.3 \times 10^{-5}$; for the compactly supported Mat\'ern
cluster kernel it gives $b = R$ exactly, with no tail margin required; for nested
Thomas the buffers of the two levels add, giving $b = 4\sigma_1 + 4\sigma_2$,
which over-covers the tight bound $4\sqrt{\sigma_1^2 + \sigma_2^2}$ and is
therefore conservative. Points are never clipped to the boundary; the only
boundary operation is the containment filter at step 5.

% -----------------------------------------------------------------------------
\subsection{Validation}
\label{ssec:verification}
% NOTE: this is a CONTRIBUTION, not an implementation note. The exact
% finite-window Var[N(W)] and the 175% discrepancy with the canonical identity
% deserve prominence -- probably a figure showing the gap as a function of nu.
%
% The $g$ and $K$ definitions have been MOVED HERE from the old Background,
% because this is the only place in the paper that uses them.
% -----------------------------------------------------------------------------

% --- moved definitions -------------------------------------------------------
% [MOVED FROM: old Section 2.2 -- trim to only what is used below]
Throughout this section, let $X$ denote a point process on $\R^d$. We write
$u, v \in \R^d$ for points (precise locations) and $du$, $dv$ for infinitesimal
regions of volume $|du|$, $|dv|$ centered at $u$ and $v$, respectively.

    \begin{definition}[Second-order product density]
        The second-order product density
        $\rho^{(2)}: \R^d \times \R^d \rightarrow \R^+$
        is defined such that $\rho^{(2)}(u, v)\,du\,dv$ is the probability of
        jointly observing a point of $X$ in each of the infinitesimal regions
        $du$ and $dv$.
    \end{definition}

    \begin{definition}[Pair correlation function]
        The pair correlation function $g(u, v): R^d \times R^d \rightarrow \R^+$
        is defined
        \[
        g(u, v) = \frac{\rho^{(2)}(u, v)}{\rho(u)\rho(v)}.
        \]
    It is clear that for any Poisson process we have $g \equiv 1$. Intuitively,
    $g$ can suggest whether SPPs are attractive ($g > 1$), or repulsive
    ($g < 1$).
    \end{definition}

    \begin{definition}[Stationarity \& isotropy]
        A point process $X$ is said to be:
            \begin{itemize}
                \item \emph{Stationary} if $X \sim X + s$ for every
                $s \in \R^d$ where $X + s = \{x + s: x \in X\}$.
                \item \emph{Isotropic} if $X \sim RX$ for every rotation $R$
                about the origin.
            \end{itemize}
    \end{definition}

    \begin{note}
        Stationarity implies constant intensity $\rho(u) = \lambda$, but the
        converse is not necessarily the case. If $X$ is stationary and isotropic,
        $g(u, v)$ depends on $u$ and $v$ only through their inter-point distance
        $r = \|u - v\|$. Abusing notation, we may write $g(r)$ for this common
        value. We now define the $K$ function.
    \end{note}

    \begin{definition}[$K$-function]
        For stationary and isotropic $X$, the $K$-function
        $K(r): \R^+ \rightarrow R^+$ is defined
            \[
            K(r) = \int_{b_r(0)} g(\|h\|) dh
            \]
        Intuitively, $\rho(u)K(r)$ is the expected number of points we expect to
        see within a distance $r$ of a typical point in $X$.
    \end{definition}

\todo{consider double duty for these definitions: they also give the cleanest
statement of what a radially averaged second-order statistic CANNOT see, which
is the argument Section 3.2 needs. Cross-reference rather than duplicate.}

% --- validation proper -------------------------------------------------------
% [MOVED FROM: old Section 3.1.3 -- unchanged]
Due to every downstream result being dependent on a diligent cloud-generation
process, we validate it against the established closed-form second-order theory.
On $W = [0,1]^2$ we compare $[[3000]]$ realizations at each of three parameter
triples spanning the design, chosen to span the overlap index: a tight-cluster
regime ($\nu = 0.07$), a moderate regime ($\nu = 0.22$), and a near-Poisson
regime ($\nu = 1.00$).
\todo{fix the $[[3000]]$ placeholder}

For the first moment, $\mathbb{E}[N(W)] = \kappa\mu|W|$. For the second, the
canonical identity $\mathrm{Var}[N(W)] = \kappa\mu(1+\mu)|W|$ holds only in the
limit of a window large relative to $\sigma$, and is badly biased in precisely
the near-Poisson regime we wish to stress. Applying the law of total variance to
$N(W) = \sum_p N_p(W)$ with
$N_p(W) \mid p \sim \mathrm{Poisson}(\mu q(p))$, $q(p) = \int_W k(x-p)\,dx$, and
evaluating the shot-noise term by Campbell's formula gives the exact
finite-window result
    \begin{equation}
    \label{eq:varN}
    \mathrm{Var}[N(W)] = \kappa\mu + \frac{\kappa\mu^2}{4\pi\sigma^2}J(\sigma)^2,
    \qquad
    J(\sigma) = \int_{-1}^{1}(1-|t|)\,e^{-t^2/4\sigma^2}\,dt,
    \end{equation}
for $W = [0,1]^2$, where the reduction to a one-dimensional integral follows from
the triangular density of the difference of two independent
$\mathrm{Uniform}(0,1)$ variables. As $\sigma \to 0$,
$J(\sigma) \to 2\sigma\sqrt{\pi}$ and \eqref{eq:varN} collapses to
$\kappa\mu(1+\mu)$, confirming the two agree in the appropriate limit. The gap
between them reaches $175\%$ at $\nu = 1$, so \eqref{eq:varN} is used as the
comparison target.

\todo{FIGURE: canonical identity vs \eqref{eq:varN} as a function of $\nu$, with
the empirical points overlaid. This makes the 175\% claim visible in one glance
and is the cheapest high-value figure in the paper.}

\begin{table}[h]
\centering
\begin{tabular}{lccccc}
\toprule
Case & $\nu$ & $\mathbb{E}[N]$ theory & empirical
     & $\mathrm{Var}[N]$ theory & empirical \\
\midrule
A & 0.07 & 200.0 & $200.4 \pm 0.6$ & 982.0 & $972.5 \pm 25.1$ \\
B & 0.22 & 100.0 & $100.2 \pm 0.4$ & 545.2 & $544.8 \pm 14.1$ \\
C & 1.00 & 20.0  & $20.0 \pm 0.1$  & 43.6  & $46.0 \pm 1.2$ \\
\bottomrule
\end{tabular}
\caption{Simulator validation against \eqref{eq:varN}. All deviations are within
$2$ standard errors.}
\label{tab:sim-validation}
\end{table}

We additionally verify that the empirical pair correlation function and Ripley's
$K$ match
    \[
    g(r) = 1 + \frac{e^{-r^2/4\sigma^2}}{4\pi\kappa\sigma^2},
    \qquad
    K(r) = \pi r^2 + \frac{1 - e^{-r^2/4\sigma^2}}{\kappa},
    \]
that the marginal offspring intensity on $W$ shows no boundary deficit
(edge-to-interior density ratio $1.003$--$1.048$ across the three cases, in all
cases consistent with unity and never in the direction an insufficient buffer
would produce), and that the displacement kernel is isotropic with per-axis
standard deviation matching $\sigma$. Estimates of $g$ normalized per realization
carry a negative bias that grows with $\mathbb{E}[N^2]$ and is therefore larger
for clustered patterns than for a matched-intensity Poisson control; we account
for this when reading residual deviations rather than attributing them to the
sampler.

% -----------------------------------------------------------------------------
\subsection{Splits and evaluation protocol}
% NOTE: RENAMED from "Splits". This is the evaluation protocol for the whole
% paper and was previously invisible, buried at the bottom of data generation.
% Section 7.3 should reference this subsection directly.
% -----------------------------------------------------------------------------
% [MOVED FROM: old Section 3.1.4]
Data are split by parameter vector, never by realization, so that no parameter
value appears in both training and evaluation. Beyond the random test split we
generate two further evaluation sets: a grid over the design grid, permitting
error to be reported as a surface over $(\log\nu, \log\lambda)$ rather than as a
single scalar; and a set drawn one increment outside each range, characterizing
degradation under extrapolation.

\todo{add here: the classical baselines (minimum contrast, Palm likelihood) will
be run on the SAME clouds. State that now -- it is what makes Section 7.3 a
comparison rather than a table of unanchored numbers.}
\fix{The classical baselines introduced in Section 2 are run on the same clouds,
under the same train/test/grid/extrapolation splits rather than on separate data
or numbers taken from the literature. Section 7.3 is therefore a comparison
against a shared  test set.}

\todo{state the metric explicitly (logMAE per parameter?) -- the \texttt{\textbackslash lm}
macro existed in the old preamble but was never used in the text}
\fix{Throughout Section 7 we report mean squared error in log-normalized parameter
space. The scalar figures in Section 7.3 average this quantity over the test
set and over the three parameters; where relevant we report the per-parameter
loss breakdown. Because the standardization is the one used for training, 
a loss of $\mathcal{L}=1$ corresponds to an error of one training-set standard 
deviation of $\log\theta$.}

% NOTE: the "design distribution acts as a prior" remark that used to close this
% subsection has been MOVED to Section 8 (Discussion), where it belongs as a
% limitation rather than as a passing aside.


% #############################################################################
% 6. METHOD: PARAMNET
% #############################################################################
% NOTE -- PURPOSE: describe the pipeline. Present it honestly as an instance of
% a known pattern (cf. Section 2), with the genuinely interesting choices
% flagged as CHOICES to be evaluated in Section 7 -- not as conclusions.
%
% GLOBAL EDIT REQUIRED IN THIS SECTION:
% every sentence of the form "we tested many variants and present the best"
% must become "we compare X, Y, Z; results in Section 7.x". There are at least
% three such sentences and they are what currently deletes the comparison
% thread from the paper.
% #############################################################################
\section{Method: ParamNet}
\label{sec:method}

% [MOVED FROM: old Section 3 opening]
Let $\pc$ denote a point cloud, $v$ its vectorized persistence summary, and
$e \in \R^{128}$ the resulting feature vector. ParamNet follows the pipeline
    \begin{equation}
    \label{eq:pipeline}
    \pc \xrightarrow{\;\text{Vectorization}\;} v
    \xrightarrow{\;\text{Encoder}\;} e
    \xrightarrow{\;\text{FCNN}\;} \hat\theta.
    \end{equation}

\todo{add one sentence here conceding that this pipeline shape follows the
neural-estimator pattern of related work [A], and that the contribution lies in
the feature stage. Better to say it plainly here than to have a reviewer say it.}
\fix{We note that the pipeline shape follows the neural-estimator patterns used in REF
and REF. Our contribution lies in the feature-extraction stage, the application to
a wider range of SPP families, and the careful evaluation of design choices.}

The model is trained and evaluated on synthetic clouds generated by sweeping a
design distribution of the parameter space. The synthetic clouds are processed
together before training so that the feature-extraction step is performed only
once.

Feature extraction is the model's key component: it is where the cloud's
topological summary is vectorized and made usable for machine learning. Much
work has been done to develop vectorizations that ensure stability and minimize
the inevitable information loss. We explore several such methods here, hoping
that their relative efficacies provide insight into the structure they capture.
\todo{THIS SENTENCE IS THE ORIGINAL RESEARCH OBJECTIVE and it is the best
statement of it anywhere in the draft. Promote the idea to the Introduction's
contributions list and point forward to Section 7.1 from here.}
\fix{Feature extraction is the model's key component: it is where the cloud's 
topological summary is vectorized and made usable for machine learning. 
A substantial body of work has developed vectorizations of persistence 
diagrams —- persistence images, Betti curves, landscapes, kernel methods -- 
that trade off stability, resolution, and information loss in different ways, 
but there is little consensus on which is preferable for recovering physical 
generating parameters from a single observed point cloud. We implement and 
compare several such vectorizations against this task directly, rather than 
picking one on prior grounds; the comparison, and what it reveals about the 
structure each vectorization does and does not preserve, is the subject of 
Section~\ref{ssec:exp-summaries}.}


The vectorized summary $v$ is passed through an encoder to obtain the
feature vector $e \in \R^{128}$, which is then given to the FCNN
$NN_\phi: \R^{128} \to \R^{n_params}$ to produce the estimate
$\hat\theta = (\hat\kappa, \hat\mu, \hat\sigma)$.

\todo{$NN_\phi$ is never actually specified -- depth, widths, activations.}
\fix{$NN_\phi$ is a small feed-forward network: two hidden layers of width 64 and 
32, each followed by a ReLU nonlinearity and dropout ($p=0.1$), then a final 
linear layer mapping to the $n_{params}=3$ outputs $(\hat\kappa, \hat\mu, \hat\sigma)$.}

% -----------------------------------------------------------------------------
\subsection{Filtration and vectorization}
% -----------------------------------------------------------------------------
% [MOVED FROM: old Section 3.2]
Although we implemented and tested a variety of vectorized topological summaries,
we focus primarily on the most effective: persistence image.
\todo{REWRITE -- this sentence deletes the comparative contribution in half a
line. Replace with: "we implement and compare N vectorizations; the comparison
is the subject of Section 7.1. For the remainder of this section we describe the
pipeline instantiated with persistence images."}
We expose the pipeline as follows;
    \begin{enumerate}
        \item Pick filtration function and construct filtrations
        $\{(\filt)_{r\geq 0}\}_i$.
        \item Compute persistence diagrams $\{D_0\}_i$, $\{D_1\}_i$
        (see Definition \ref{def:per_image})
        \item Calibrate diagrams to determine common birth/persistence axis
        ranges, shared across the dataset, then vectorize each diagram into a
        fixed-size raster using those ranges.
    \end{enumerate}
The key design choices that follow are:
    \begin{itemize}
        \item Filtration function
        \item Calibration procedure
        \item Image resolution
    \end{itemize}
Although the latter two points can be optimized with a grid-search, the choice of
filtration function plays a fundamental role in the optimal extraction of
information. Various exploratory experiments suggested that the canonical Rips
filtration does not recover sufficient information about the local density of the
clouds, and therefore does not provide the model with a complete picture.
\todo{"various exploratory experiments suggested" is an unsupported assertion
carrying a genuinely interesting empirical finding. Either promote the evidence
to Section 7.1, or state it here as a hypothesis and forward-reference.}

In the present iteration of ParamNet, we use the $DTM_k$ filtration, as seen in
Example \ref{ex:dtm}, which provides a richer summary. Furthermore, the parameter
$k$ allows us to compute summaries at various scales, which may be combined to
provide the model with a more complete picture.

The simultaneous use multiple $k$-scales per point cloud introduces complexities
in the pipeline - specifically how the vectorized persistence summaries are
encoded and combined to produce the final feature vector for $NN_\phi$. We
discuss this in the following section.

% -----------------------------------------------------------------------------
\subsection{Calibration}
\label{ssec:calibration}
% NOTE: this is CONTRIBUTION #2 and it is currently three levels deep. The
% material itself is in good shape -- the fix is positional, plus a figure.
% Old cross-references pointed at \ref{sec:calibration}, which never existed;
% label is now \ref{ssec:calibration}.
% -----------------------------------------------------------------------------
% [MOVED FROM: old Section 3.2.1 -- unchanged]
Consider the synthetic training clouds $\{\pc_i, \theta_i\}_{i=1}^{N_{train}}$
and fix the homology dimension $h \in \{0,1\}$ for some choice of filtration. We
obtain $N_{train}$ persistence diagrams $\{D^{(h)}_i\}_{i=1}^{N_{train}}$. Each
diagram will have a different distribution of points, some of which may have
extreme outliers with respect to the full set of diagrams.

In order to ensure that in the resulting persistence images each pixel bears the
same meaning, we need to calibrate the diagrams to have the same axis ranges. The
naive approach of calibrating to the maximum deaths and births in order to
include all points is inadequate due to the resulting relative positioning of the
bulk of points.

If we have a diagram with very strong outliers and chose universal axes that
include these outliers, the remaining diagrams will be rescaled in a way that
places the majority of their points near the origin. The diagrams -- and
therefore their vectorizations -- will lose a lot of texture.

\todo{FIGURE: naive min/max calibration vs coverage-quantile calibration, same
diagram, side by side. The "loss of texture" argument is visual and is currently
carried entirely by prose.}

Our solution is to calibrate according to \emph{per-diagram coverage quantiles}.
For each training diagram $D_i^{(h)}$, define the three summary statistics
\begin{equation}
    b_i^{\min} = \min_{(b,d) \in D_i^{(h)}} b, \qquad
    b_i^{\max} = \max_{(b,d) \in D_i^{(h)}} b, \qquad
    p_i^{\max} = \max_{(b,d) \in D_i^{(h)}} (d - b),
\end{equation}
i.e. the smallest birth, largest birth, and largest persistence value attained
\emph{within a single diagram}. Given a coverage level $q \in (0,1)$ and a
padding factor $\kappa \geq 1$, we set the calibrated birth and persistence
ranges to be the $q$-quantiles of these per-diagram statistics, taken over the
training set:
\begin{equation}
    \big[\,Q_{1-q}\big(\{b_i^{\min}\}_{i=1}^{N_{train}}\big),\
    \kappa \cdot Q_{q}\big(\{b_i^{\max}\}_{i=1}^{N_{train}}\big)\,\big]
    \qquad \text{and} \qquad
    \big[\,0,\ \kappa \cdot
    Q_{q}\big(\{p_i^{\max}\}_{i=1}^{N_{train}}\big)\,\big],
\end{equation}
where $Q_q(\cdot)$ denotes the empirical $q$-quantile. In practice we take
$q = 0.99$ and $\kappa = 1.05$, which were determined via grid-search. These two
ranges are shared by every diagram of homology dimension $h$, for every point
cloud in the dataset (train, validation, test, and adversarial alike), so that a
given pixel of the resulting persistence image always refers to the same region
of birth-persistence space.
\todo{NOTATION CLASH: $\kappa$ is the padding factor here and the parent
intensity everywhere else in the paper. Rename the padding factor.}

The key difference with the naive min/max calibration is \emph{what} is being
quantiled. Pooling all points from all diagrams and taking a quantile over that
pooled set implicitly weights each diagram by how many topological features it
produces: a single diagram with an unusually large number of points -- near the
diagonal or otherwise -- can dominate the resulting bound, while a diagram
consisting of one strong outlier and few other points contributes only that
outlier. By instead computing one representative statistic per diagram and
quantiling \emph{over diagrams}, every training cloud contributes exactly once to
the calibration, regardless of how many persistence pairs it happens to produce.
Framing outlier rejection as a coverage fraction over diagrams -- rather than
over points -- makes the trade-off explicit and controllable: at most
$\lceil (1-q) N_{train} \rceil$ diagrams may have their most extreme birth or
persistence value fall outside the calibrated frame (and are clipped rather than
discarded), while the bulk of the population retains full resolution. The naive
approach of calibrating to the global maximum is recovered as the degenerate case
$q \to 1$.

This scheme is also compatible with the stability guarantees of persistence
images \citep{adams2017persistence}, which are stated for a fixed bandwidth and a
fixed image domain: since $\sigma_{\text{pixels}}$ is specified in pixel units
(Definition \ref{def:per_image}) rather than in absolute birth/persistence units,
and the birth and persistence ranges are calibrated independently of one another,
a pixel remains a comparable, well-defined location across every image built from
the calibrated imager -- irrespective of the very different physical spans the
two axes typically have.

\paragraph{Train/validation/test leakage.}
Because the calibrated ranges are a property of the training distribution, they
must be fit \emph{exclusively} on $\{D_i^{(h)}\}_{i \in \text{train}}$, then
frozen and applied unchanged to validation, test, and adversarial diagrams.
Fitting the ranges on the full, unsplit population of clouds -- before the
train/validation/test partition is drawn -- lets information about the held-out
clouds leak into the calibration used to build every training image, since every
pixel's meaning is defined relative to the calibrated bounds. As the
train/validation/test partition is redrawn per random seed, calibration must
likewise be refit per seed, on that seed's training rows only, and then applied
frozen to the corresponding validation, test, and adversarial rows.

\paragraph{Degenerate axes.}
For homology dimension $H_0$, birth values are often identically zero (or
otherwise carry negligible spread) under the chosen filtration, in which case
$b_i^{\min} \approx b_i^{\max}$ for all $i$ and the calibrated birth range
collapses to a point. We detect this case explicitly and fall back to a
one-dimensional vectorization of the diagram rather than fabricate an axis, since
a degenerate range would make an entire column of pixels convey no information.

% -----------------------------------------------------------------------------
\subsection{Encoder and fusion}
% NOTE: the ARCHITECTURE description stays here; the SWEEP RESULTS that used to
% be interleaved with it have been moved to Section 7.2. Keep the split clean:
% this subsection says what the options are, Section 7.2 says which won.
% -----------------------------------------------------------------------------
% [MOVED FROM: old Section 3.3]
Many variations of encoders were tested for this model. We will focus on the
variation that produced the best results, noting that the encoding step is an
important aspect that will require further study.
\todo{REWRITE -- same problem as the vectorization subsection. State the design
space here; report the outcome in Section 7.2.}

Consider our training clouds $\{\pc_i, \theta_i\}_{i=1}^{N_{train}}$. We pick $K$
scales for our $DTM_k$ filtration, and compute persistence images for $D$
homology dimensions with resolution $H \times W$. Therefore for batch size $B$,
the objective is to pick a suitable encoder architecture for the following
transformation:
    \begin{equation}
    (B, K, D, H, W) \xrightarrow{\;\text{Encoder}\;} (B, 128).
    \end{equation}

When it comes to using multiple images per homology dimension for each cloud, a
central question is whether each vectorization should go through a shared-weight
encoder, or be assigned its own independent encoder.

The shared encoder is designed as follows: each $(D, H, W)$ persistence-image
stack is first augmented with two coordinate channels -- fixed $x$- and
$y$-coordinate grids linearly spaced over $[-1, 1]$ -- following the CoordConv
construction \citep{liu2018coordconv}, so that the convolutional stack has access
to absolute pixel position rather than having to infer it.

This matters here specifically because the birth and persistence axes are
calibrated (Section~\ref{ssec:calibration}): a fixed pixel location has a fixed,
shared meaning across the dataset, and CoordConv lets the network exploit that
directly instead of relearning it.

The coordinate-augmented $(D+2, H, W)$ tensor is then passed through a small
convolutional stack: three $3\times3$ convolution blocks with channel widths
$(32, 64, 128)$, each followed by a ReLU, and by $2\times2$ max-pooling with
spatial dropout between blocks (but not after the last one); the resulting
feature map is reduced by adaptive average pooling to a single spatial location
and projected by a linear layer to the embedding dimension $E$.

To apply this one encoder to all $K$ scales, the $k$-axis is folded into the
batch dimension for a single forward pass -- one call over $B \cdot K$ images
rather than $K$ separate calls -- so the weights are tied across scales by
construction rather than by any explicit sharing mechanism.

This produces a $(B, K, E)$ tensor. The next question concerns how the per-$k$
summary is prepared for $NN_\phi$. The options explored include:
    \begin{enumerate}
        \item Flat concatenation of the $K$ per-scale embeddings into a single
        $(B, K \cdot E)$ vector
        \item A small 1-D convolutional fusion over the ordered $k$-axis (two
        $\text{Conv1d}$ layers, letting adjacent scales mix), followed by
        average-pooling over $k$ into a fixed-width $(B, F)$ vector whose size
        does not depend on $K$
        \item The same convolutional fusion, but flattening rather than pooling
        over $k$, giving a $(B, F \cdot K)$ vector that preserves which position
        came from which scale at the cost of a width that again grows with $K$
    \end{enumerate}

\todo{end this subsection by forward-referencing Section 7.2 rather than
announcing the winner here}


% #############################################################################
% 7. EXPERIMENTS
% #############################################################################
% NOTE -- PURPOSE: this section is the structural fix for the whole paper.
% It does three jobs, in this order:
%   7.1 the ORIGINAL research question (which topological summary is best) --
%       currently absent from the draft, must be reconstructed from the
%       exploratory work already done;
%   7.2 architecture ablations -- ALREADY DONE, just relocated here;
%   7.3 estimation performance -- THE EMPTY SLOT. Headers written now so the
%       report reads as planned rather than unfinished.
%
% Open the section with a short paragraph stating the experimental protocol
% (pointing at Section 5.4) and the metric, so it is not repeated three times.
% #############################################################################
\section{Experiments}
\label{sec:experiments}

\todo{opening paragraph: protocol (Section \ref{sec:simulator}), metric, seeds,
hardware, and how the three subsections relate}

% -----------------------------------------------------------------------------
\subsection{Comparison of topological summaries}
\label{ssec:exp-summaries}
% NOTE: THE RE-ELEVATED ORIGINAL OBJECTIVE. This subsection is the paper's
% clearest claim to a contribution independent of the "best model" thread, and
% it currently exists nowhere in the draft -- only as three deleted sentences
% ("we tested many, here is the best").
%
% WHAT GOES HERE:
%   * filtration function: Rips vs DTM_k, and DTM across k. This is where the
%     "Rips does not recover local density" claim from Section 6.1 finally gets
%     evidence.
%   * vectorization: persistence image vs Betti curve vs whatever else was run.
%   * ideally crossed, or at least one held fixed while the other varies.
%   * homology dimension: H_0 alone, H_1 alone, both.
%   * INTERPRETATION, not just a table -- which structural information does each
%     summary retain, and which parameter does it help most? A per-parameter
%     breakdown here would make the section genuinely interesting rather than a
%     ranking.
% -----------------------------------------------------------------------------
\todo{WRITE THIS SUBSECTION. Much of the underlying work is already done; it
needs to be recovered from the exploratory experiments and run to a consistent
protocol (same seeds, same splits, same design grid).}

\todo{TABLE: filtration $\times$ vectorization, test loss mean $\pm$ sd over
seeds}

\todo{FIGURE: per-parameter error by summary type -- expect $\sigma$ and
$\kappa$ to behave very differently}

% -----------------------------------------------------------------------------
\subsection{Architecture ablations}
\label{ssec:exp-ablations}
% NOTE: THIS IS ALREADY DONE -- it was hiding as parenthetical prose inside the
% encoder description. Six configurations x five seeds plus a ten-seed
% confirmation, with a real finding (instability, not just worse averages).
% Presented as results it is a legitimate experimental subsection.
% -----------------------------------------------------------------------------
% [MOVED FROM: old Section 3.3 -- encoder-mode sweep]
Experiments suggested that a shared encoder works just as well as the independent
option: we ran a sweep crossing \emph{encoder mode} (shared-weight vs.
independent-per-$k$) against \emph{fusion mode}, five seeds per combination, on
the same $DTM_{5,10,15}$ nested-Thomas setup used throughout. Restricting to the
simplest fusion strategy (flat concatenation), the shared-weight encoder matched
the independent one to within seed noise:

\begin{center}
\begin{tabular}{lccc}
\toprule
encoder mode & fusion & test loss (mean $\pm$ sd) & failed seeds \\
\midrule
shared      & concat & $0.200 \pm 0.009$ & $0/5$ \\
independent & concat & $0.217 \pm 0.028$ & $0/5$ \\
\bottomrule
\end{tabular}
\end{center}
\todo{promote to a proper \texttt{table} environment with a caption and label;
report the full six-configuration cross, not just the two concat rows}

with a ten-seed run of the shared-weight configuration ($0.205 \pm 0.012$)
confirming the same order of magnitude. Since tying weights across scales also
cuts the parameter count roughly $K$-fold and removes one design axis, we adopt
the shared encoder going forward.

% [MOVED FROM: old Section 3.3 -- fusion sweep]
We determined once again that the simplest option was the most effective, and we
concatenate the per-$k$ outputs. The full sweep, crossing both encoder mode and
fusion strategy (six combinations, five seeds each), showed that the two
convolutional-fusion variants are not simply worse on average but
\emph{unstable}: individual seeds collapse to a test loss of $\sim\!0.9$--$1.0$,
while other seeds of the very same configuration land right on the concatenation
baseline ($\sim\!0.20$).

Only the independent-encoder, flatten-pooled combination converged cleanly on all
five seeds, but even so it did not outperform plain concatenation
($0.212 \pm 0.005$ vs.\ $0.200 \pm 0.009$). Flat concatenation was therefore
preferred not only for its simplicity, but because it was the only strategy --
together with its shared/independent-encoder counterpart -- that trained reliably
across seeds without any sign of the collapse observed for the
convolutional-fusion alternatives.

\todo{FIGURE: per-seed test loss scatter across the six configurations. The
instability finding is the interesting part and a table of means hides it.}

% -----------------------------------------------------------------------------
\subsection{Estimation performance}
\label{ssec:exp-performance}
% #############################################################################
% *** THE EMPTY SLOT ***
% This subsection does not exist yet and is the main outstanding work item.
% Writing the headers now is what makes the progression report read as PLANNED
% rather than UNFINISHED. Do not delete these TODOs -- they are the plan.
% #############################################################################
% -----------------------------------------------------------------------------

\paragraph{Comparison against classical estimators.}
\todo{minimum contrast on $K$ and on $g$, and Palm likelihood, run on the SAME
clouds (Section \ref{sec:simulator}). This is the headline comparison and the
only one that can substantiate the limitations claimed in Section 1.}

\paragraph{Error surfaces over the design grid.}
\todo{error as a surface over $(\log\nu, \log\lambda)$ using the grid evaluation
set. Expect degradation as $\nu \to 1$ (weak identifiability) -- if the model
degrades more gracefully than minimum contrast there, that IS the result.}

\paragraph{Per-parameter breakdown.}
\todo{$\kappa$, $\mu$, $\sigma$ separately. They are not equally hard and a
single scalar loss hides this.}

\paragraph{Performance relative to the irreducible error floor.}
\todo{report error as a ratio to the single-realization floor defined in
Section \ref{sec:problem}, so the numbers mean something absolute.}

\paragraph{Extrapolation.}
\todo{the out-of-range evaluation set -- degradation one increment outside each
design range.}

\paragraph{Cross-family results.}
\todo{Thomas / Mat\'ern cluster / nested Thomas. Whether one model trained across
families beats per-family models is an interesting question in itself, and the
common reparameterization of Section 5.1 is what makes it askable.}


% #############################################################################
% 8. DISCUSSION AND LIMITATIONS
% #############################################################################
% NOTE -- PURPOSE: pre-empt the obvious objections. Short.
% #############################################################################
\section{Discussion and Limitations}

\paragraph{The design distribution is a prior.}
% [MOVED FROM: end of old Section 3.1.4, where it was a stray closing remark]
Finally, we note that the design distribution acts as a prior, so the trained
model approximates a posterior summary under it, and its ranges must therefore
cover the intended application domain.
\todo{expand -- this is the single most important limitation of any amortized
estimator and deserves a full paragraph, not a trailing sentence}

\paragraph{Computational cost.}
\todo{persistence computation scales badly in $n$; this is why the point budget
of Section 5.1 is capped at 800. State the cost explicitly and compare against
the cost of minimum contrast.}

\paragraph{Single realization.}
\todo{the estimator sees one realization; classical methods often assume the
same, but say so explicitly}

\paragraph{Homogeneity and stationarity.}
\todo{everything here assumes homogeneous, stationary, isotropic processes}


% #############################################################################
% 9. FUTURE WORK
% #############################################################################
% NOTE -- PURPOSE: REQUIRED for the progression report; for the paper this can
% be folded into Section 8. Keep it concrete and sequenced -- this is the
% section the panel reads to judge whether there is a viable thesis here.
% #############################################################################
\section{Future Work}

\todo{immediate: complete Section \ref{ssec:exp-performance}. Give this a
timeline -- it is the nearest-term deliverable.}

\todo{near term: the encoder question flagged in Section 6.3 as requiring
further study -- attention over scales, set-based encoders operating on diagrams
directly rather than on rasters}

\todo{medium term: inhomogeneous / non-stationary processes, where classical
closed forms are least available and the topological approach has the most to
gain -- this is arguably where the thesis-scale contribution lies}

\todo{medium term: real data. Galaxy catalogues if the running example of
Section 1 is cosmological; name a specific dataset.}

\todo{longer term: uncertainty quantification -- the model currently emits a
point estimate with no interval, which limits its use as a drop-in replacement
for likelihood-based methods}


% #############################################################################
% BIBLIOGRAPHY
% #############################################################################
\todo{set up references.bib -- currently only \texttt{anai2020dtm},
\texttt{adams2017persistence}, and \texttt{liu2018coordconv} are cited, and the
Diggle / Tanaka references in Section 1 are inline text rather than citations}

\bibliographystyle{plainnat}
\bibliography{references}

\end{document}